% Figure: a toggle (over-centre) clamp near its locked position. A hand force
% on the handle drives two links whose common pin approaches the line joining
% their outer pins; as the pin crosses that line the clamping force runs up
% toward the singular value and the mechanism locks. Panel (a) the linkage;
% panel (b) the FBD of the middle pin as a three-force node.
\begin{figure}[htbp]
\centering
\begin{tikzpicture}[
  force/.style={draw=engblue, -{Stealth}, very thick},
  rxn/.style={draw=engred, -{Stealth}, very thick},
  link/.style={draw=engblack, line width=2.2pt},
  scale=1.0,
]
% ----- panel (a): the two-link toggle -----
\begin{scope}[shift={(0,0)}]
\coordinate (A) at (0,0);      % fixed pin (clamp arm side)
\coordinate (B) at (4.0,0);    % fixed pin (handle pivot side)
\coordinate (K) at (2.0,0.55); % knee pin, just above line AB
% base line AB (dashed reference)
\draw[draw=enggray, dashed, thin] (A) -- (B) node[midway,below=1pt,font=\scriptsize,enggray] {toggle line};
% two links
\draw[link] (A) -- (K);
\draw[link] (B) -- (K);
% pins
\foreach \p in {A,B,K} {\draw[draw=engblack, fill=enggray!20, thick] (\p) circle (0.10);}
\node[font=\scriptsize, anchor=north east] at (A) {$A$};
\node[font=\scriptsize, anchor=north west] at (B) {$B$};
\node[font=\scriptsize, anchor=south] at (K) {$K$};
% handle continues past B
\draw[link] (B) -- ++(1.5,0.9);
\draw[force] ($(B)+(1.5,0.9)$) -- ++(0.35,0.75) node[anchor=south west, font=\small, engblue] {$F_h$};
\node[font=\scriptsize, enggray, anchor=west] at ($(B)+(1.5,0.9)+(0.1,-0.2)$) {handle};
% clamp output at A (the arm presses the workpiece down)
\draw[force] (A) -- ++(0,-1.0) node[anchor=north, font=\small, engblue] {$Q$};
\node[font=\scriptsize, anchor=north east] at (1.9,-0.9) {(a)};
% small angle epsilon at knee
\node[font=\scriptsize, enggray, anchor=south] at (2.0,0.68) {$\varepsilon$};
\end{scope}

% ----- panel (b): FBD of the knee pin K -----
\begin{scope}[shift={(7.6,0.1)}]
\coordinate (Kp) at (0,0);
\fill[engblack] (Kp) circle (0.05);
\node[font=\scriptsize, anchor=south] at (Kp) {$K$};
% link forces along each link (nearly collinear, small angle to horizontal)
\draw[rxn] (Kp) -- ++(-2.3,-0.6) node[anchor=east, font=\small, engred] {$F_{AK}$};
\draw[rxn] (Kp) -- ++(2.3,-0.6) node[anchor=west, font=\small, engred] {$F_{BK}$};
% driving force from the handle, pushing the knee down toward the toggle line
\draw[force] (Kp) -- ++(0,1.3) node[anchor=south, font=\small, engblue] {$D$};
% angle markers
\draw[draw=enggray, thin] (-0.9,-0.235) arc (180+14.6:180:0.9);
\node[font=\scriptsize, enggray, anchor=east] at (-0.95,-0.13) {$\varepsilon$};
\node[font=\scriptsize, anchor=north] at (0,-0.95) {(b)};
\end{scope}
\end{tikzpicture}
\caption[A toggle (over-centre) clamp near lock]{An over-centre toggle clamp
approaching its locked position. Panel (a): two links join the fixed pins $A$
and $B$ at a knee pin $K$ that sits a small distance above the straight line
$AB$. A hand force $F_h$ on the handle drives the knee down toward that line,
and the clamp arm at $A$ presses the work with the output force $Q$. Panel
(b): the free-body diagram of the knee pin $K$, a three-force node carrying the
driving force $D$ from the handle and the two nearly collinear link forces
$F_{AK}$ and $F_{BK}$, each inclined at the small angle $\varepsilon$ to the
toggle line. Force balance at the knee gives $F_{AK}=F_{BK}=D/(2\sin
\varepsilon)$, which runs away as $\varepsilon\to 0$: the clamping force grows
without bound as the knee nears the toggle line, so a modest hand force
develops a large clamp force at the instant of lock. Driving the knee just past
the line puts the mechanism over centre, where the work's reaction can no
longer push it back, and the clamp holds with the hand removed.}
\label{fig:vol03ch01:toggle-clamp}
\end{figure}
